\section{Introduction}
A classifier trained with noisy annotations must learn useful visual structure while receiving incorrect supervision. Its final class probabilities summarize the selected model, but omit how its support for an annotation changed during learning. Training dynamics expose this information for examples encountered during optimization: dataset cartography describes confidence and variability, and margin-based ranking identifies potentially mislabeled examples \citep{swayamdipta2020dataset,pleiss2020identifying}. A deployment image, however, has no measured training trajectory.

We investigate a direct way to make part of this signal available at inference. A class-conditioned regressor learns an early-history summary from frozen image features and an observed annotation. At inference, the regressor takes the base model's predicted class as its conditioning label. Its output estimates the margin summary associated with that image--class pair. A monotone transformation of the output provides a reliability ranking and weights predictions across augmented views.

The central empirical question is whether the historical target contributes beyond learning a predictor of a static checkpoint statistic. We therefore keep the regressor architecture, feature dimension, training budget, and downstream weighting rule fixed when comparing an early-history mean, the mean over all recorded epochs, and a static margin. We separately compare features of the selected base classifier with an independent frozen ResNet18 encoder. This design distinguishes the target comparison from the effect of the feature source.

The evaluated early-history variant improves accuracy on both CIFAR-10N Worse and CIFAR-100N, with gains repeated on three independently trained CIFAR-100N bases and three CIFAR-10N Worse bases. The early summary outperforms the full-history mean in these exploratory CIFAR runs. The feature-source comparison also favors features from the selected base over those of the independent encoder for the early target. We report calibration and an ensemble baseline alongside these gains: an ensemble of two CIFAR-100N bases reaches higher accuracy than either base with the proposed weighting.

Our contribution is an empirical study of when early-training information
helps at inference under noisy supervision. We compare class-conditioned
scalar weighting with matched targets and evaluate posterior compression
in two prespecified external tests. CIFAR-10.1 measures residual
distillation against the original classifier. CIFAR-10.2 tests direct
distillation from five early checkpoints against one checkpoint with five
views. Neither test meets its prespecified primary accuracy criterion.
The latter experiment finds a teacher-level advantage that does not carry
over to the corresponding one-pass students. Independently trained bases
are the replication unit. Regressor initializations are nested within
each base.

